\begin{partdivider}{V}{QUELLE SOCIÉTÉ APRÈS L'IA ?}\end{partdivider}

\bookchap{14}{Trois futurs possibles}

Au fil de ce livre, nous avons exploré les forces de transformation à l'œuvre : l'abondance naissante de l'intelligence, la recomposition du travail, la crise de la vérité, la concentration du pouvoir infrastructurel et la mutation des organisations. Aucune de ces forces ne pointe vers une direction unique. Ensemble, elles dessinent un espace de possibles, un éventail de futurs entre lesquels nos choix collectifs devront trancher.

La prospective n'a pas pour vocation de prédire ; elle a pour fonction de rendre les choix lisibles. Pour éclairer les décisions qui nous attendent, il est utile de construire des scénarios contrastés. Non pas des prédictions, mais des \emph{possibles} exagérés pour mieux en saisir la logique interne.

Construisons donc trois futurs, en apparence extrêmes, qui représentent autant de trajectoires cohérentes selon la manière dont nous répondrons à la question de la rareté, du pouvoir et du sens.

\begin{scenariobox}{A}{L'IA comme multiplicateur humain}
Dans ce premier scénario, l'IA accomplit la promesse la plus optimiste de ses partisans : elle devient le «bicycle de l'esprit» dont rêvait Steve Jobs, un multiplicateur universel des capacités humaines.

La condition de ce scénario est double. D'une part, l'accès à l'IA est réellement démocratisé : les modèles ouverts dominent, les infrastructures de calcul sont régulées comme des services publics ou des biens communs, et chaque citoyen dispose d'un droit d'accès de base. D'autre part, le système éducatif a réussi sa mutation : il forme massivement des esprits capables de questionner, de vérifier et de superviser des systèmes intelligents.

Les conséquences sont profondes. La productivité individuelle explose. L'entrepreneuriat se démocratise : n'importe qui, muni d'une idée et d'une suite d'agents, peut monter un projet viable. Les coûts de production des biens cognitifs s'effondrent, faisant baisser le prix de nombreux services (santé, éducation, conseil, design), ce qui augmente le pouvoir d'achat réel des populations.

Le travail ne disparaît pas, mais il est réinventé. Les tâches d'exécution étant automatisées, l'humain se consacre à la définition des objectifs, à la créativité, à la relation et à l'éthique. De nouvelles professions émergent, centrées sur l'orchestration de systèmes hybrides et l'audit des machines. L'expertise n'est plus la possession d'un savoir statique, mais l'agilité à naviguer dans un océan d'intelligence synthétique.

Dans ce scénario, la société est plus fluide, plus innovante, potentiellement plus égalitaire si les gains de productivité sont redistribués. La tension principale n'est plus l'accès à la capacité cognitive, mais la préservation du lien social face à une hyper-productivité individualiste. Le risque est celui de l'isolement : entouré d'agents sur mesure, l'individu pourrait se couper de la confrontation avec la pensée divergente et l'altérité humaine, s'enfermant dans des bulles cognitives ultra-personnalisés.
\end{scenariobox}

\begin{scenariobox}{B}{L'IA comme nouvelle concentration du pouvoir}
Le deuxième scénario incarne la matérialisation de nos craintes les plus vives concernant l'oligopolisation et le néo-féodalisme technologique.

Dans ce scénario, l'hypothèse de l'abondance d'intelligence pour l'utilisateur final n'est qu'une façade. En réalité, l'essentiel de la valeur et du pouvoir est capté par une minorité : les propriétaires des modèles fondationnels, des centres de données et des données d'entraînement. Le \emph{compute} est la nouvelle terre, et ceux qui la possèdent perçoivent la rente.

L'IA n'est pas un bien commun, c'est un service sous location. Les entreprises et les individus dépendent de quelques API propriétaires pour penser, produire et décider. Le coût de l'abonnement aux systèmes de pointe augmente avec leur puissance, créant une asymétrie croissante entre ceux qui peuvent se payer l'IA de pointe et les autres.

La société se stratifie en trois couches. Au sommet, les propriétaires de l'infrastructure et les architectes de systèmes, qui orientent l'économie et le discours. Au milieu, une classe de travailleurs augmentés, hautement productifs mais précaires, engagés dans une compétition féroce (l'effet Reine rouge) pour accomplir les tâches de supervision et d'orchestration que la machine ne peut pas encore faire seuls. En bas, une majorité de la population, reléguée à des tâches physiques non automatisables (soins, services, logistique), ou à des micro-tâches d'entraînement ou de validation de l'IA (le travail du clic), pour des salaires misérables.

La productivité globale a augmenté de manière spectaculaire, mais les bénéfices ont été accaparés par le sommet. Les institutions démocratiques sont érodées : la capacité de fabriquer la réalité à grande échelle permet aux plateformes de modeler l'opinion, tandis que la complexité technique des systèmes rend toute régulation politique inopérante ou obsolète.

Ce scénario n'est pas une fiction. Il est la prolongation de tendances déjà observées avec les plateformes du Web 2.0. Si l'IA suit la même logique d'intégration verticale et d'effet de réseau sans contrepoids institutionnel, il est le plus probable.
\end{scenariobox}

\begin{scenariobox}{C}{La société hybride}
Le troisième scénario refuse tant l'utopie de la multiplication harmonieuse que la dystopie de la concentration. Il postule que l'avenir sera un pragmatisme incertain : une société hybride, où humains et IA coexistent, coopèrent et s'ajustent mutuellement par tâtonnements.

Dans ce scénario, l'IA n'est ni un bien commun parfait, ni un outil de domination total. Elle est une infrastructure mixte. Des modèles ouverts coexistent avec des services propriétaires de pointe. La régulation fonctionne par à-coups, souvent en retard d'une guerre, mais elle parvient à briser les concentrations les plus dangereuses, à imposer de la transparence sur les algorithmes de recommandation et à garantir des droits fondamentaux numériques.

Le travail est en profonde mutation, mais la transition est gérée — non sans douleur — par des filets de sécurité sociale, des politiques de formation continue et des expérimentations de réduction du temps de travail. L'automatisation détruit des emplois, mais la société invente de nouvelles formes de contribution, reconnues et rémunérées : les métiers du care, de la réparation, de la médiation, de l'artisanat physique et de l'accompagnement humain retrouvent une valeur sociale et économique centrale.

L'école a bifurqué. Elle n'est plus le temple de la restitution du savoir, mais le lieu où l'on apprend la coopération humain-machine, la pensée critique et la navigation dans l'incertitude. Les compétences essentielles dans ce scénario ne sont plus la rapidité d'exécution ou la mémoire, mais :

\begin{eplist}
\item \textbf{Le jugement} : savoir quand et pourquoi désobéir à la machine ou s'en méfier.
\item \textbf{La responsabilité} : assumer les conséquences des décisions, même déléguées.
\item \textbf{La créativité de situation} : inventer des solutions dans des contextes que la machine n'a jamais vus.
\item \textbf{La confiance} : certifier l'authenticité humaine et bâtir des relations fiables dans un monde de synthèse.
\end{eplist}
Ce scénario est le moins spectaculaire, mais peut-être le plus résilient. Il ne suppose pas que la technologie soit intrinsèquement bonne ou mauvaise, mais il exige une vigilance démocratique constante, une capacité d'adaptation institutionnelle et un refus de laisser le marché seul dicter l'ordre du monde. Il suppose aussi que nous redéfinissions ce qui donne de la valeur à une vie humaine.
\end{scenariobox}

Une interprétation s'impose : aucun de ces trois scénarios n'est «le bon». Ils ne sont pas mutuellement exclusifs ; des éléments de l'un apparaîtront dans les autres. Leur fonction est de nous forcer à préciser ce que nous voulons.

Si nous voulons éviter le Scénario B, il faut des politiques volontaristes de rupture des monopoles et de soutien à l'open source. Si nous voulons tendre vers le Scénario A, il faut une révolution de l'éducation et de l'accès. Si nous visons le pragmatisme du Scénario C, il faut réinventer notre État-providence et nos mécanismes de solidarité.

Car au cœur de tous ces scénarios se trouve une question qui dépasse la technologie : celle de la valeur. Si l'intelligence cesse d'être rare, sur quel fondement allons-nous organiser notre vie ensemble ?

\bookchap{15}{Et si nous devions réinventer la société ?}

Depuis l'aube de la révolution industrielle, nos sociétés se sont organisées autour d'un principe simple, rarement formulé mais profondément intériorisé : \textbf{la valeur économique découle de la productivité, la productivité découle du travail, et le travail découle de l'effort humain.} Ce triptyque a structuré notre rapport au temps, à la dignité, à la redistribution et au pouvoir. L'IA vient frapper ce triptyque d'une incongruité historique. Pour la première fois, la productivité peut découler de l'effort machine, et non plus de l'effort humain. Si une machine peut produire, analyser et décider (dans certaines limites) à un coût marginal proche de zéro, le lien entre productivité et travail humain se trouve distendu.

Cela ne signifie pas que le travail humain deviendra inutile. Nous avons vu qu'il se transforme et se déplace. Mais cela signifie que la part de la richesse produite grâce au seul travail humain va probablement diminuer, tandis que la part produite grâce au capital technologique (l'IA) va augmenter.

\begin{aiquote}
Comment redistribuer la valeur dans une économie où la création de richesse dépend de plus en plus du capital cognitif et de moins en moins de l'effort humain ?
\end{aiquote}

C'est ici que l'enjeu n'est plus technologique, mais profondément politique et moral.

Première question : \textbf{la redistribution.} Si l'IA augmente la productivité globale tout en réduisant la demande de travail humain pour certaines tâches, les gains de productivité ne couleront pas naturellement vers les salariés. Ils auront tendance à être captés par les propriétaires du capital (les entreprises qui déploient l'IA) et par les propriétaires de l'infrastructure (les fournisseurs de modèles et de compute). Sans mécanisme de redistribution, l'IA pourrait être la plus grande machine à creuser les inégalités jamais inventée.

Les pistes sont connues, mais difficiles à mettre en œuvre. On parle de revenu universel, de taxation des robots, de dividendes computationnels (où chaque citoyen toucherait une part des revenus générés par les modèles entraînés sur les données publiques), ou encore de réduction massive du temps de travail sans perte de revenu, pour partager le travail restant. Aucune de ces pistes n'est parfaite ; toutes supposent un État capable de taxer le capital immatériel hautement mobile, ce qui est un défi majeur pour les finances publiques.

Deuxième question : \textbf{le sens du travail.} Nous avons tendance à confondre travail et emploi. L'emploi est une construction économique ; le travail est une activité humaine qui donne sens, structure le temps, forge l'identité et tisse le lien social. Si l'emploi tel que nous le connaissons se contracte ou mute, le besoin de travail — au sens d'activité créatrice et utile — ne disparaît pas. Il faudra peut-être élargir notre définition de ce qui compte comme une contribution valuable. Aujourd'hui, notre système ne rémunère que les activités qui ont un marché. Mais de nombreuses activités vitales pour la société — élever des enfants, prendre soin des personnes âgées, entretenir le lien de voisinage, produire des biens communs open source, s'engager dans la vie associative — ne sont pas ou peu rémunérées. Si l'IA libère du temps en automatisant l'exécution, ce temps libéré pourrait être réinvesti dans ces activités de care et de lien, à condition que la société choisisse de les reconnaître et de les financer.

Troisième question : \textbf{la mesure de la valeur.} Le Produit Intérieur Brut (PIB) mesure la valeur des échanges monétaires. Si l'IA permet de produire gratuitement des biens qui étaient autrefois payants (traductions, logiciels, analyses, contenus éducatifs), le PIB pourrait stagner ou baisser, alors même que le bien-être réel et l'accès aux services augmentent. Inversement, des externalités négatives majeures (la fatigue informationnelle, la pollution cognitive, l'érosion de la confiance) ne sont pas captées par le PIB. Nous aurons besoin de nouveaux indicateurs de richesse et de santé sociale pour naviguer dans ce monde nouveau.

Quatrième question : \textbf{l'autonomie.} Nous avons évoqué à plusieurs reprises le risque de dépendance cognitive. Si l'IA pense avec nous, il est probable qu'une part de notre intuition et de notre mémoire s'atrophie, comme notre sens de l'orientation s'est atrophié avec le GPS. Jusqu'où sommes-nous prêts à déléguer ? Y a-t-il un noyau dur de l'agir humain que nous refusons d'externaliser, non pas parce que la machine le ferait mal, mais parce que \emph{le faire} est constitutif de notre humanité ?

Penser par soi-même, décider en son âme et conscience, créer de ses propres mains, affronter l'adversité, se tromper et apprendre de ses erreurs : tout cela a une valeur qui n'est pas économique, mais existentielle. Une société qui déléguerait l'intégralité de ses choix et de sa pensée à des machines, fût-ce pour son plus grand confort matériel, ne serait pas une société humaine.

\begin{aiquote}
Comment préserver l'autonomie et la liberté dans un monde où l'IA peut penser, choisir et agir à notre place ?
\end{aiquote}

La réponse ne réside pas dans le refus de la technologie, mais dans la délibération éthique. Nous devons cartographier les zones de délégation acceptable (calculer un itinéraire, synthétiser un rapport) et les zones de souveraineté non négociable (le jugement moral, l'éducation des enfants, le vote, le soin des vulnérables). Cinquième question : \textbf{la gouvernance.} Aucune de ces questions — redistribution, sens, mesure, autonomie — ne trouvera de réponse dans le marché seul. Le marché suit la productivité ; il ne se soucie ni de la justice ni de l'autonomie. Ces réponses exigent des espaces de délibération politique légitimes. Or, nos institutions démocratiques, conçues pour des nations industrielles du XX\up{e} siècle, semblent souvent dépassées par la vitesse et la technicité du changement.

Nous aurons besoin d'une nouvelle ingénierie démocratique. Des assemblées citoyennes sur l'IA, des autorités de régulation algorithmiques dotées de moyens d'audit, des traités internationaux sur les données et la souveraineté numérique, une éducation permanente des élus et des fonctionnaires. L'enjeu n'est pas de freiner l'IA, mais d'accélérer la capacité de la société à la comprendre, à la critiquer et à l'orienter.

Ce que nous retiendrons de ce chapitre : Si l'hypothèse de ce livre est juste — si l'intelligence devient abondante —, alors nous ne pouvons pas nous contenter de greffer l'IA sur nos anciennes institutions. Nous devons réinventer les règles du jeu. Non pas parce que l'IA le commande, mais parce que les possibilités qu'elle ouvre nous confrontent à nos propres responsabilités.

La question n'est plus «Que va-t-il se passer ?», mais «Que choisissons-nous de faire advenir ?».
